Analog radio-over-fiber links at millimeter-wave frequencies require equalization methods that address component nonlinearities while supporting practical receiver implementation. This work investigates a framework for online neural equalization of a single-carrier 16-QAM optical-beating link near 60 GHz. A modular physical model represents electrical drive, Mach--Zehnder modulation, fiber propagation, photodetection, wireless transmission, and receiver recovery. Independent training, validation, and test sequences support comparisons between a least-mean-square feed-forward equalizer and neural receivers with configurable hidden-layer dimensions and activation functions. The neural formulation uses continuous in-phase and quadrature outputs, mean-square-error training, and chronological sample-by-sample stochastic gradient descent. All stored neural quantities use signed Q8.24 arithmetic, with explicit rounding and saturation boundaries and backward propagation based on pre-update weights. The evaluation protocol links bit-error rate, error-vector magnitude, and convergence behavior to the resource, throughput, and energy requirements of a candidate FPGA implementation. Initial training sequences and periodic supervised bursts provide the intended mechanism for adaptation, with the first hardware phase operating at a controlled sample rate. The framework is designed to identify operating conditions in which nonlinear neural compensation provides a useful advantage and to assess the implementation cost of maintaining that advantage during online training.
